% Generated draft: requires mathematical and source review.
\documentclass[11pt]{article}
\usepackage[margin=25mm]{geometry}
\usepackage{fontspec}
\setmainfont{Latin Modern Roman}
\usepackage{amsmath,amssymb,mathtools,unicode-math}
\setmathfont{Latin Modern Math}
\usepackage{enumitem,needspace,xurl,hyperref}
\setlength{\parindent}{0pt}
\setlength{\parskip}{5pt}
\setlength{\emergencystretch}{3em}
\newcommand{\N}{\mathbb{N}}
\newcommand{\Z}{\mathbb{Z}}
\newcommand{\Q}{\mathbb{Q}}
\newcommand{\R}{\mathbb{R}}
\newcommand{\C}{\mathbb{C}}
\newcommand{\F}{\mathbb{F}}
\newcommand{\A}{\mathbb{A}}
\newcommand{\PP}{\mathbb{P}}
\newcommand{\E}{\mathbb{E}}
\newcommand{\eps}{\varepsilon}
\newcommand{\dd}{\,\mathrm{d}}
\begin{document}
\section*{UnsolvedMath ID: 30004941}
\textbf{Source problem code:} catalogue entry
\subsection{Definitions and mathematical statement}
The supplied record concerns the amplituhedron \(A_{n,k,m}(Z)\), but does not define it or specify the admissible parameters. The following is a conventional candidate definition, included conditionally rather than asserted to be the exact convention of the source. Let \(n,k,m\) be nonnegative integers with \(k+m\leq n\). Write \(\operatorname{Gr}(k,n)\) for the Grassmannian of \(k\)-dimensional linear subspaces of \(\R^n\), and let \(\operatorname{Gr}_{\geq0}(k,n)\) consist of the row spaces of rank-\(k\) matrices \(C\in\R^{k\times n}\) whose ordered maximal minors are all nonnegative. In a column convention for the external data, let \(Z\in\R^{(k+m)\times n}\) have every ordered maximal minor strictly positive. The conventional amplituhedron is
\[
 A_{n,k,m}(Z)=\bigl\{\operatorname{rowspan}(CZ^{\mathsf T}):\operatorname{rowspan}(C)\in\operatorname{Gr}_{\geq0}(k,n)\bigr\}\subseteq\operatorname{Gr}(k,k+m).
\]
The intended parameter range, any parity restrictions on \(m\), and the source's precise positivity and matrix conventions require confirmation.

To explain the flow terminology, let \(M\) be a smooth manifold containing the amplituhedron, let \(g\) be a Riemannian metric on \(M\), and let \(f:M\to\R\) be smooth. A negative gradient flow satisfies
\[
 \frac{\dd}{\dd t}\Phi_t(x)=-\operatorname{grad}_g f\bigl(\Phi_t(x)\bigr),\qquad \Phi_0(x)=x.
\]
Forward preservation of \(A=A_{n,k,m}(Z)\) means that the flow is defined on \(A\) for every \(t\geq0\) and
\[
 \Phi_t(A)\subseteq A\qquad(t\geq0).
\]
This definition is an explicit possible interpretation: the source does not specify the ambient manifold, metric, potential, sign convention, or whether preservation is forward or two-sided.

\emph{Requested classification and existence problem.} For fixed admissible parameters \((n,k,m)\), classify the gradient flows preserving \(A_{n,k,m}(Z)\). Determine whether, for every admissible positive \(Z\), there exists such a flow with additional properties sufficient to establish a homeomorphism
\[
 A_{n,k,m}(Z)\cong\overline B^{\,d},\qquad
 \overline B^{\,d}=\left\{x\in\R^d:\sum_{i=1}^{d}x_i^2\leq1\right\},
\]
for the intended nonnegative integer dimension \(d\). The record does not specify \(d\), the equivalence relation for the classification, or the additional dynamical hypotheses needed for this topological conclusion. Preservation alone is insufficient: the constant potential gives the identity flow, which preserves every subset without imposing closed-ball topology. Consequently, this is a conditional statement schema requiring clarification, rather than a fully specified conjecture.
\subsection{Short English statement}
Classify gradient flows preserving the amplituhedron \(A_{n,k,m}(Z)\), and determine whether every admissible positive \(Z\) admits a suitable flow establishing that the amplituhedron is homeomorphic to a closed ball. The supplied record leaves the admissible parameters, precise flow class, classification equivalence, and sufficient dynamical conditions unspecified.
\subsection{Sources}
\begin{itemize}
\item Supplied local UnsolvedMath registry record, permanent ID \(30004941\), title \emph{Gradient Flows and Ball Topology of Amplituhedra}, source problem code \texttt{catalogue entry}, in the Oberwolfach Reports Open Problems collection. No individual proposer is supplied.
\item Supplied primary-source citation: \emph{Mini-Workshop: Scattering Amplitudes, Cluster Algebras, and Positive Geometries} (2022), Oberwolfach Reports, pp.\ 3101--3133, DOI \texttt{10.4171/owr/2021/57}. The extracted problem identifies p.\ 3127. \url{https://doi.org/10.4171/owr/2021/57}. The registry separately records proposed year 2021 and the tag \texttt{owr-year:2021}; the supplied record does not reconcile these with the citation's 2022 date.
\item Supplied literature reference: Karp--Williams, \emph{Gradient Flows and Adjoint Orbits}. \url{https://doi.org/10.1007/s00220-022-04540-5}.
\item Supplied literature reference labeled ``original Oberwolfach report'': \url{https://doi.org/10.4171/owr/2021/57}.
\item Literature assessment checked at \texttt{2026-08-22}: the supplied dated curation labels the problem open and states that gradient-flow methods establish closed-ball topology for several important amplituhedron families, while the requested classification and suitable-flow existence for every positive \(Z\) remain open. This is supplied dated curation, not an independently verified literature assessment; its status and cited claims require independent verification.
\item The registry labels the statement \texttt{corrected\_verified} and reports direct checking of the official report. These are supplied provenance assertions; no independent source visit or mathematical verification is claimed here.
\end{itemize}
\end{document}
